\subsection{Password lifecycle in the stable design}
\label{sec:password}
Only the Creator workspace exposes controls for setting or changing an itinerary password. The agent sets the password for a trip and sends the protected export to the intended travelers. The agent can communicate the password separately from the file. A traveler may enter the password to open a received package. The Viewer has no control for setting the password, changing it, or publishing a replacement export.

A password change is a publishing operation. The Creator retains the trip identifier, reserves a higher revision, generates a fresh salt and nonce, and encrypts a new snapshot using the new password. This applies even when the schedule is identical. The Creator may then distribute the new file through the same channel as the original. Password protection belongs to each exported file. Earlier files still use their original passwords. A traveler can keep reading the schedule already saved on the device.

At import, the Viewer reads the format and checks bit 2 of the native envelope's flags to decide whether to prompt. An import cancelled at the prompt leaves the library as it was. The import path checks framing and length before prompting. The flag becomes authenticated only after the GCM tag verifies. For a protected file, the Viewer holds the supplied password for key derivation and uses the salt and nonce carried in that file. If authentication fails, the Viewer reports the unlock failure and leaves the existing itinerary intact. After successful authentication, the Viewer decompresses and validates the snapshot, checks the trip ID and revision, and stores a new trip or replaces the matching saved version. A stale or equal revision does not replace the saved schedule. Replacement preserves the checklist and completion flag.

The import password is needed again for a later protected package, even when that package updates a trip already on the phone. For a password-only revision, the traveler enters the new password at import. The old password will not unlock the newly encrypted file. Once the new version is saved, the traveler reads it from local storage without an itinerary-password prompt. The Viewer discards the transfer password and derived key after use. Operating-system and storage controls govern access to the extracted files.

For example, an agent issues Kashmir revision 3 with password A, then changes the password and sends revision 4 with password B. A traveler who already imported revision 3 can continue reading it. The Viewer asks for B when opening revision 4. After validating the package, it updates the one locally saved Kashmir trip. Subsequent access to that saved trip needs neither A nor B. The example describes the expected behavior of the stable design.

The shared password grants access to anyone who learns it. A short PIN remains guessable offline despite the cost of Argon2id. The Creator's exclusive editing controls belong to the application interface. To establish that a revision came from the original agent, the Viewer would need a trusted publisher key and a signature covering the trip identifier, revision, and package contents. The standard mode derives its key from a common seed that can be recovered from the software.
